% ===============================
% Worksheet / Manual Template
% ===============================
\documentclass[12pt, a4paper]{article}

% ===============================
% Metadata
% ===============================
\newcommand*{\CourseCode}{MAT321}
\newcommand*{\CourseTitle}{Real Analysis II}
\newcommand*{\Chapter}{Subspaces, Separable Spaces}
\newcommand*{\Topics}{Dense Set and Perfect Set}
\newcommand*{\DocDate}{February 08, 2026}

\include{header}

% ==========================================
% Document
% ==========================================
\begin{document}
\FrontPage
\Large


\begin{heading}
    Subspaces
\end{heading}

\vspace{10pt}

\begin{definition}[Subspace and Induced Metric]
Let $(X,d)$ be a metric space and let $Y$ be a non-empty subset of $X$.
The restriction of the metric $d$ to $Y\times Y$, denoted by $d_Y$, is a metric on $Y$
called the \emph{induced metric}. The metric space $(Y,d_Y)$ is called a
\emph{subspace} of $(X,d)$.
\end{definition}

\vspace{100pt}

\begin{problem}[Examples of Subspaces]
Verify that:
\begin{enumerate}
    \item the closed unit interval $[0,1]$ and the set of all rational numbers $\mathbb{Q}$
    are subspaces of $\mathbb{R}$ with the usual metric;
    \item the unit circle, the closed unit disc, and the open unit disc are subspaces of
    the metric space $(\mathbb{C},d)$ of complex numbers;
    \item the real line $\mathbb{R}$ itself is a subspace of the complex plane $\mathbb{C}$.
\end{enumerate}
\end{problem}

\clearpage

\begin{definition}[Open Sphere in a Subspace]
Let $(X,d)$ be a metric space, let $Y\subseteq X$, and let $y\in Y$.
The open sphere in the subspace $(Y,d_Y)$ centered at $y$ with radius $r>0$ is denoted by
$S_r^{\,Y}(y)$ and is defined by
\[
S_r^{\,Y}(y)=\{x\in Y : d_Y(x,y)<r\}.
\]
\end{definition}

\vspace{100pt}

\begin{problem}[Relation Between Open Spheres]
Show that for every $y\in Y$ and every $r>0$,
\[
S_r^{\,Y}(y)=S_r(y)\cap Y,
\]
where $S_r(y)$ denotes the open sphere in $(X,d)$.
\end{problem}


\clearpage

\begin{theorem}[Open Sets in Subspaces]
Let $(X,d)$ be a metric space and let $Y\subseteq X$. A subset $A$ of $Y$ is open in
$(Y,d_Y)$ if and only if there exists a set $G$ open in $(X,d)$ such that
\[
A = G \cap Y.
\]
\end{theorem}

\clearpage

\begin{theorem}[Open Sets in a Subspace vs Ambient Space]
Show that if $Y\subseteq X$ and a subset of $Y$ is open in $(X,d)$, then it is also open in
the subspace $(Y,d_Y)$. Show by examples that the converse need not be true.
\end{theorem}

\vspace{20pt}

\begin{problem}[Open in a Subspace may Not be open in Superspace]
Let $X=\mathbb{R}$ with the usual metric and let
\[
Y=[0,1].
\]
Then the open sphere in $(Y,d_Y)$ centered at $0$ with radius $\tfrac12$ is
\[
S^{\,Y}_{1/2}(0)=[0,\tfrac12),
\]
which is open in $Y$ but not open in $\mathbb{R}$. Note that $Y$ itself is not open in
$\mathbb{R}$.
\end{problem}

\vspace{80pt}

\begin{problem}[Example: Real Line vs Complex Plane]
Show that an open interval of the real line is not an open subset of the complex plane
$\mathbb{C}$ with the usual metric.
\end{problem}

\clearpage

\begin{heading}
    Dense Sets
\end{heading}

\vspace{10pt}

\begin{definition}[Dense Set]
A subset $A$ of a metric space $(X,d)$ is said to be \emph{dense} (or \emph{everywhere dense})
in $X$ if the closure of $A$ is $X$, that is,
\[
\overline{A}=X.
\]
\end{definition}

\clearpage

\begin{heading}
    Nowhere Dense Sets
\end{heading}

\vspace{10pt}

\begin{definition}[Nowhere Dense Set]
Let $A$ be a subset of a metric space $(X,d)$. The set $A$ is said to be
\emph{nowhere dense} in $X$ if the interior of its closure is empty, that is,
\[
\operatorname{int}(\overline{A})=\emptyset.
\]
\end{definition}

\vspace{10pt}

\begin{problem}[Basic Properties of Nowhere Dense Sets]
Show that:
\begin{enumerate}
    \item every finite subset of $X$ is nowhere dense;
    \item every subset of a nowhere dense set is nowhere dense;
    \item $A$ is nowhere dense in $X$ if and only if $\overline{A}$ is nowhere dense in $X$.
\end{enumerate}
\end{problem}

\clearpage

\begin{heading}
    Somewhere Dense Sets and Dense-in-itself Sets
\end{heading}

vspace{10pt}

\begin{definition}[Somewhere Dense Set]
A subset $A$ of $X$ is said to be \emph{somewhere dense} in $X$ if it is not nowhere dense.
Equivalently, $A$ is somewhere dense in $X$ if and only if $\overline{A}$ contains a non-empty
open sphere.
\end{definition}

\vspace{10pt}

\begin{definition}[Dense-in-itself Set]
A subset $A$ of a metric space $(X,d)$ is said to be \emph{dense-in-itself} if every point of
$A$ is a limit point of $A$, that is,
\[
A\subseteq A'.
\]
\end{definition}

\clearpage

\begin{heading}
    Perfect Sets
\end{heading}

vspace{10pt}

\begin{definition}[Perfect Set]
A subset $A$ of a metric space $(X,d)$ is said to be \emph{perfect} if it is closed and
dense-in-itself, that is,
\[
A=A'.
\]
\end{definition}

vspace{10pt}

\begin{problem}[Examples of Perfect Sets]
Show that every closed interval, the empty set $\emptyset$, and the whole space $X$ are
perfect sets.
\end{problem}

\clearpage

\begin{problem}[Cantor Set is Perfect]
Show that the Cantor set is a perfect set.
\end{problem}

\clearpage

\begin{heading}
    Separable Metric Space
\end{heading}

\begin{definition}[Separable Metric Space]
A metric space $(X,d)$ is said to be \emph{separable} if there exists a countable subset of $X$
which is dense in $X$.
\end{definition}

\clearpage

\begin{problem}[Euclidean Space is Separable]
Show that the Euclidean space $\mathbb{R}^n$ is separable.
\end{problem}

\clearpage

\begin{problem}[$\ell^p$ is Separable]
Show that the metric space $\ell^p$ of all sequences $(x_n)$ of real numbers such that $\sum_{n=1}^\infty |x_n|^p<\infty$ with the metric
\[d_p((x_n),(y_n))=\left(\sum_{n=1}^\infty |x_n-y_n|^p\right)^{1/p}\]
is separable for every $1\leq p<\infty$.
\end{problem}

\clearpage

\begin{problem}[$\ell^\infty$ is Not Separable]
Prove that the metric space $\ell^\infty$ of all bounded sequences with the supremum metric
is not separable.
\end{problem}


\clearpage

\EndPage
\end{document}